%!TEX root = TQFTnotes.tex
%%%%%%%%%%%%%%%%%%%%%%%%%%%%%%
\chapter{3-2-1 TQFT and modular categories}\label{c:321}

References: for modular categories, \ocite{etingof-fusion}*{};
for modular functor, \ocite{bakalov-kirillov} and works by
Jurgen Andersen (FIXME). For correspondence between 321 theories
and modular categories, \ocite{bdspv}.

\section{3-2-1 theories}\label{s:321-theories}
Overview of 3-2-1 theory. Theorem:
every 3-2-1 theory defines on $\calC$ a structure of
ribbon category.

\section{Graphical calculus}\label{s:321-graphical}

\section{Modular group action and modular categories}\label{s:modular-group-action}
Verlinde algebra. Links.
Theorem: a 3-2-1 theory defines on $\calC$ structure of a modular category
(assuming it is semisimple).

\section{Reshetikhin--Turaev theory}\label{s:rt-theory}
Theorem: every modular category defines a (central extension of) 3-2-1 theory.
Historical reference; surgery.

\section{Chern--Simons theory}\label{s:chern-simons}
Category $\calC(\g, k)$. Result of Kazhdan--Lusztig--Finkelberg. Relation with
path integral (Axelrod--Singer)

\section{Link invariants}\label{s:link-invariants}
Invariants of links in $S^3$. Jones polynomial.